% main_probability_metropolis.tex
\documentclass[aspectratio=169]{beamer}

% --- Language & encoding ---
\usepackage[utf8]{inputenc}
\usepackage[T1]{fontenc}
\usepackage[english,french]{babel}
\usepackage{lmodern}
\usepackage{microtype}

% --- Math & tables ---
\usepackage{amsmath,amssymb}
\usepackage{booktabs}
\usepackage{graphicx}
\usepackage{changepage}

% --- Metropolis Theme ---
\usetheme[numbering=fraction,progressbar=frametitle]{metropolis}
\setbeamertemplate{frame numbering}[fraction]

% --- Colors (same palette as prior lectures) ---
\definecolor{BootcampBlueDark}{HTML}{1A3C68}
\definecolor{TextBlack}{HTML}{000000}
\definecolor{AccentGold}{HTML}{DAA520}
\definecolor{Crimson}{HTML}{990000}
\definecolor{MatrixGreen}{HTML}{228B22}
\definecolor{MatrixRed}{HTML}{B22222}

\setbeamercolor{block title example}{bg=BootcampBlueDark!90!black,fg=white}
\setbeamercolor{block body example}{bg=BootcampBlueDark!5!white,fg=black}

% === Thick frame title band ===
\setbeamercolor{frametitle}{bg=BootcampBlueDark,fg=white}
\setbeamerfont{frametitle}{series=\bfseries,size=\large}
\setbeamertemplate{frametitle}{
  \nointerlineskip
  \begin{beamercolorbox}[wd=\paperwidth,ht=3.2ex,dp=1.4ex,leftskip=1em,rightskip=1em]{frametitle}
    \usebeamerfont{frametitle}\insertframetitle
  \end{beamercolorbox}
}

% === Custom Definition Block ===
\newenvironment{definitionblock}[1]{%
  \par\vspace{0.4em}%
  \begin{beamercolorbox}[
      wd=\textwidth,
      sep=0.4em,
      leftskip=0.3em,
      rightskip=0.6em,
      rounded=false,
      shadow=false
    ]{block title example}%
    \raisebox{0.15em}{\usebeamerfont{block title example}#1}%
  \end{beamercolorbox}%
  \vspace{-0.4em}%
  \begin{beamercolorbox}[
      wd=\textwidth,
      sep=1em,
      leftskip=0.6em,
      rightskip=0.6em,
      rounded=false,
      shadow=false
    ]{block body example}%
    \usebeamerfont{block body example}%
}{%
  \end{beamercolorbox}%
  \par\vspace{0.6em}%
}

% === Global text formatting ===
\setbeamercolor{normal text}{fg=TextBlack,bg=white}
\setbeamercolor{title}{fg=TextBlack}
\setbeamercolor{structure}{fg=BootcampBlueDark}
\setbeamercolor{progress bar}{fg=BootcampBlueDark,bg=gray!30}

% --- Course Info ---
\title{Probability (Lecture 3) --- Foundations \& Random Variables}
\author{Your Name}
\institute{Your Program / Department}
\date{\today}
% \titlegraphic{\includegraphics[height=1.2cm]{your_logo.pdf}}

\metroset{
  sectionpage=progressbar,
  subsectionpage=progressbar
}

% --- Section & subsection transitions (plain) ---
\AtBeginSection{
  \frame[plain]{\sectionpage}
}
\AtBeginSubsection{
  \frame[plain]{\subsectionpage}
}

\begin{document}

\maketitle

% ===== Outline =====
\begin{frame}{Outline}
  \tableofcontents[hideallsubsections]
\end{frame}

% ===========================
% SECTION 0 — Motivation
% ===========================
\section{Motivation \& Overview}

\begin{frame}{Outline}
  \tableofcontents[currentsection, hideothersubsections]
\end{frame}

\subsection{Why Probability in Finance?}
\begin{frame}{Uncertainty, Modeling, \& Decisions}
\begin{itemize}
  \item Model uncertainty in asset prices, defaults, and risk factors.
  \item Support portfolio analysis, pricing, and risk measurement.
  \item Link to optimization via expectations, variance, and covariances.
\end{itemize}
\end{frame}

\subsection{Learning Goals}
\begin{frame}{What You Will Learn}
\begin{definitionblock}{Learning Goals}
By the end of this lecture, you will be able to:
\begin{itemize}
  \item Define probability spaces and manipulate probabilities.
  \item Apply conditional probability, independence, total probability, and Bayes' rule.
  \item Work with random variables: pmf/pdf/cdf, expectation, variance, covariance.
  \item Understand joint/marginal distributions and the (multi)variate normal.
\end{itemize}
\end{definitionblock}
\end{frame}

% ======================================
% SECTION 1 — Basics
% ======================================
\section{Basics of Probability}

\begin{frame}{Outline}
  \tableofcontents[currentsection, hideothersubsections]
\end{frame}

\subsection{Probability Spaces \& Events}
\begin{frame}{Foundations}
\begin{definitionblock}{Probability Space}
A triple \((\Omega,\mathcal{F},\mathbb{P})\) where:
\begin{itemize}
  \item \(\Omega\): sample space (outcomes),
  \item \(\mathcal{F}\): \(\sigma\)-algebra of events,
  \item \(\mathbb{P}\): probability measure on \(\mathcal{F}\).
\end{itemize}
\end{definitionblock}
\begin{itemize}
  \item Events are subsets \(A\in\mathcal{F}\).
  \item Axioms: non-negativity, normalization, countable additivity.
\end{itemize}
\end{frame}

\subsection{Rules, Independence, Conditional Probability}
\begin{frame}{Rules \& Conditional Probability}
\begin{itemize}
  \item Complement: \(\mathbb{P}(A^c)=1-\mathbb{P}(A)\).
  \item Inclusion–exclusion: \(\mathbb{P}(A\cup B)=\mathbb{P}(A)+\mathbb{P}(B)-\mathbb{P}(A\cap B)\).
  \item Conditional probability: \(\mathbb{P}(A\mid B)=\dfrac{\mathbb{P}(A\cap B)}{\mathbb{P}(A)}\) if \(\mathbb{P}(B)>0\).
  \item Independence: \(\mathbb{P}(A\cap B)=\mathbb{P}(A)\mathbb{P}(B)\).
\end{itemize}
\end{frame}

\subsection{Total Probability \& Bayes}
\begin{frame}{Law of Total Probability \& Bayes' Theorem}
\begin{definitionblock}{Total Probability}
If \(\{B_i\}\) partitions \(\Omega\), then
\[
\mathbb{P}(A)=\sum_i \mathbb{P}(A\mid B_i)\,\mathbb{P}(B_i).
\]
\end{definitionblock}
\begin{definitionblock}{Bayes' Theorem}
\[
\mathbb{P}(A\mid B)=\frac{\mathbb{P}(B\mid A)\,\mathbb{P}(A)}{\mathbb{P}(B)}.
\]
\end{definitionblock}
\end{frame}

\subsection{Applications: Discrete Events in Finance}
\begin{frame}{Default Risk \& Conditional Views}
\begin{itemize}
  \item Event \(D\): default within horizon; conditional on macro state \(M\).
  \item Use total probability across scenarios; update via Bayes given signals.
  \item Joint vs conditional interpretations in portfolio/credit risk.
\end{itemize}
\end{frame}

% ======================================
% SECTION 2 — Random Variables
% ======================================
\section{Random Variables: Discrete \& Continuous}

\begin{frame}{Outline}
  \tableofcontents[currentsection, hideothersubsections]
\end{frame}

\subsection{Definition \& Examples}
\begin{frame}{What is a Random Variable?}
\begin{definitionblock}{Random Variable}
A measurable function \(X:(\Omega,\mathcal{F})\to(\mathbb{R},\mathcal{B})\).
\end{definitionblock}
\begin{itemize}
  \item Financial examples: returns, log-returns, losses, counts (defaults).
  \item Distribution of \(X\) described by pmf/pdf/cdf.
\end{itemize}
\end{frame}

\subsection{Discrete Variables}
\begin{frame}{PMF, Expectation, Variance (Discrete)}
\begin{itemize}
  \item PMF: \(p_X(x)=\mathbb{P}(X=x)\).
  \item CDF: \(F_X(x)=\sum_{t\le x} p_X(t)\).
  \item Mean: \(\mathbb{E}[X]=\sum_x x\,p_X(x)\).
  \item Variance: \(\mathrm{Var}(X)=\mathbb{E}[X^2]-\mathbb{E}[X]^2\).
\end{itemize}
% TODO: add Bernoulli/Poisson examples
\end{frame}

\subsection{Continuous Variables}
\begin{frame}{PDF, CDF, Expectation (Continuous)}
\begin{itemize}
  \item PDF: \(f_X(x)\ge 0\), \(\int_{-\infty}^\infty f_X(x)\,dx=1\).
  \item CDF: \(F_X(x)=\int_{-\infty}^x f_X(t)\,dt\).
  \item Mean: \(\mathbb{E}[X]=\int x f_X(x)\,dx\).
  \item Variance: \(\mathrm{Var}(X)=\int (x-\mu)^2 f_X(x)\,dx\).
\end{itemize}
% TODO: add Normal/Exponential quick examples
\end{frame}

\subsection{Moments \& Co-moments}
\begin{frame}{Expectation, Variance, Covariance, Correlation}
\begin{itemize}
  \item Linearity: \(\mathbb{E}[aX+bY]=a\,\mathbb{E}[X]+b\,\mathbb{E}[Y]\).
  \item Covariance: \(\mathrm{Cov}(X,Y)=\mathbb{E}[XY]-\mathbb{E}[X]\mathbb{E}[Y]\).
  \item Correlation: \(\rho_{XY}=\dfrac{\mathrm{Cov}(X,Y)}{\sigma_X\sigma_Y}\).
\end{itemize}
\[
\textcolor{BootcampBlueDark}{\text{Portfolio variance: }}\;
\mathrm{Var}(w^\top R)=w^\top \Sigma\, w.
\]
\end{frame}

\subsection{Joint \& Marginal Distributions}
\begin{frame}{Joint/Marginal \& Independence}
\begin{itemize}
  \item Joint pmf/pdf \(p_{X,Y}(x,y)\) or \(f_{X,Y}(x,y)\).
  \item Marginalization: \(f_X(x)=\int f_{X,Y}(x,y)\,dy\).
  \item Independence: \(f_{X,Y}(x,y)=f_X(x)f_Y(y)\).
\end{itemize}
% TODO: add simple bivariate example
\end{frame}

% ======================================
% SECTION 3 — Normal \& MVN
% ======================================
\section{Normal \& Multivariate Normal Distributions}

\begin{frame}{Outline}
  \tableofcontents[currentsection, hideothersubsections]
\end{frame}

\subsection{Normal Distribution}
\begin{frame}{Gaussian Basics}
\begin{itemize}
  \item \(X\sim \mathcal{N}(\mu,\sigma^2)\), pdf:
  \[
  f_X(x)=\frac{1}{\sqrt{2\pi\sigma^2}}
  \exp\!\left(-\frac{(x-\mu)^2}{2\sigma^2}\right).
  \]
  \item Standardization: \(Z=\dfrac{X-\mu}{\sigma}\).
  \item Symmetry, tail behavior, and z-scores.
\end{itemize}
\end{frame}

\subsection{Multivariate Normal}
\begin{frame}{Definition \& Properties}
\begin{itemize}
  \item \(R\sim \mathcal{N}(\mu,\Sigma)\) with mean vector \(\mu\) and covariance \(\Sigma\).
  \item Any linear combination \(w^\top R\) is normal with
    \(\mathbb{E}[w^\top R]=w^\top\mu\), \(\mathrm{Var}(w^\top R)=w^\top\Sigma w\).
  \item Conditioning formulas support regression and inference.
\end{itemize}
\end{frame}

\subsection{Finance Relevance}
\begin{frame}{Why (Multi)Normal in Finance?}
\begin{itemize}
  \item Approximate modeling of asset returns.
  \item Mean–variance portfolio analysis: inputs \((\mu,\Sigma)\).
  \item Diversification via the covariance matrix; role of correlation.
\end{itemize}
\end{frame}

% ======================================
% SECTION 4 — Applications
% ======================================
\section{Applications in Finance}

\begin{frame}{Outline}
  \tableofcontents[currentsection, hideothersubsections]
\end{frame}

\subsection{Modeling Returns \& Risk}
\begin{frame}{Asset Returns, VaR, \& Stress}
\begin{itemize}
  \item Using pdf/cdf for tail probabilities and quantiles (VaR).
  \item Portfolio risk via \(w^\top\Sigma w\); correlation and diversification.
  \item Joint distributions for co-movement and stress scenarios.
\end{itemize}
\end{frame}

\subsection{Connecting to Next Lecture}
\begin{frame}{Preview of Lecture 4}
\begin{itemize}
  \item Conditional expectation, transformations, and moment-generating tools.
  \item LLN/CLT for aggregation of risks.
  \item Applications to pricing and statistical inference.
\end{itemize}
\end{frame}

% ======================================
% SECTION 5 — Summary \& References
% ======================================
\section{Summary \& References}

\begin{frame}{Outline}
  \tableofcontents[currentsection, hideothersubsections]
\end{frame}

\subsection{Key Takeaways}
\begin{frame}{What to Remember}
\begin{itemize}
  \item Probability spaces, rules, conditional probability, independence.
  \item Random variables: pmf/pdf/cdf; expectation, variance, covariance.
  \item Joint/marginal distributions; normal and multivariate normal.
  \item Portfolio variance and diversification through covariance matrices.
\end{itemize}
\end{frame}

\subsection{References}
\begin{frame}{Suggested References}
\begin{itemize}
  \item Carl P. Simon and Lawrence Blume, \textit{Mathematics for Economists}.
  \item Sydsaeter and Hammond, \textit{Essential Mathematics for Economic Analysis}.
  \item J. Douglas Carroll and Paul E. Green, \textit{Mathematical Tools for Applied Multivariate Analysis}.
\end{itemize}
\end{frame}

\end{document}